\documentclass[12pt,a4paper]{article}
\usepackage[utf8]{inputenc}
\usepackage[spanish]{babel}
\usepackage{amsmath,amssymb}
\usepackage{geometry}
\usepackage{enumitem}
\usepackage{listings}
\usepackage{xcolor}
\usepackage{booktabs}

\geometry{margin=2cm}

\lstset{
  language=Python,
  basicstyle=\ttfamily\small,
  keywordstyle=\color{blue}\bfseries,
  stringstyle=\color{red},
  commentstyle=\color{gray}\itshape,
  showstringspaces=false,
  frame=single,
  breaklines=true,
  numbers=left,
  numberstyle=\tiny\color{gray}
}

\title{\textbf{Tecnología Digital VI --- Segundo Parcial} \\ Modelo de Examen 2}
\author{Universidad Torcuato Di Tella}
\date{}

\begin{document}
\maketitle
\noindent\textbf{Duración:} 2 horas \hfill \textbf{Nombre:} \hrulefill

\vspace{0.5cm}
\hrule
\vspace{0.5cm}

%% =============================================
\section*{Parte 1: Verdadero o Falso (2 pts c/u)}
\textit{Justifique brevemente las respuestas falsas.}

\begin{enumerate}[label=\alph*)]
    \item El Perceptrón simple puede aprender la función XOR porque tiene un peso por cada entrada más el bias.

    \item La función Softmax garantiza que las salidas sumen 1 y sean no negativas, lo que permite interpretarlas como probabilidades.

    \item En una CNN, la operación de convolución es no lineal, por lo que no es necesario agregar una función de activación después.

    \item El algoritmo ADAM combina las ideas de Momentum y Learning Rate adaptativo por parámetro.

    \item En Transfer Learning, se suelen congelar las primeras capas convolucionales porque capturan features genéricas (bordes, texturas) útiles para cualquier tarea visual.

    \item Una BiLSTM procesa la secuencia en ambas direcciones, por lo que tiene el doble de hidden states que una LSTM unidireccional.
\end{enumerate}

\vspace{0.5cm}
\hrule
\vspace{0.5cm}

%% =============================================
\section*{Parte 2: Multiple Choice (3 pts c/u)}
\textit{Marque la opción correcta. Solo una es válida.}

\begin{enumerate}[label=\arabic*)]
    \item ¿Cuál es el propósito principal de \textbf{Batch Normalization}?
    \begin{enumerate}[label=\Alph*)]
        \item Reemplazar las funciones de activación en cada capa.
        \item Normalizar las salidas de cada capa para combatir el Internal Covariate Shift y acelerar la convergencia.
        \item Reducir la cantidad de parámetros de la red.
        \item Eliminar la necesidad de usar un optimizador como ADAM.
    \end{enumerate}

    \item En Style Transfer, la \textbf{Style Loss} se calcula usando:
    \begin{enumerate}[label=\Alph*)]
        \item La diferencia pixel a pixel entre la imagen generada y la imagen de estilo.
        \item La diferencia entre las matrices de Gram de los feature maps de la imagen generada y la imagen de estilo.
        \item El error cuadrático entre los logits de la última capa FC.
        \item La divergencia KL entre las distribuciones de los feature maps.
    \end{enumerate}

    \item ¿Por qué se divide por $\sqrt{d_k}$ en el \textbf{Scaled Dot-Product Attention}?
    \begin{enumerate}[label=\Alph*)]
        \item Para que las queries y keys tengan norma unitaria.
        \item Para reducir la cantidad de parámetros del modelo.
        \item Para evitar que el producto punto crezca demasiado con la dimensión, estabilizando la softmax.
        \item Para convertir los scores en probabilidades.
    \end{enumerate}

    \item ¿Qué técnica usa \textbf{InstructGPT} para alinear las respuestas del LLM con la intención del usuario?
    \begin{enumerate}[label=\Alph*)]
        \item Data Augmentation sobre el corpus de pre-entrenamiento.
        \item Reinforcement Learning from Human Feedback (RLHF).
        \item Agregar más capas de Transformer al modelo base.
        \item Fine-tuning supervisado sin ningún componente adicional.
    \end{enumerate}

    \item En una celda LSTM, la \textbf{Forget Gate} se encarga de:
    \begin{enumerate}[label=\Alph*)]
        \item Decidir qué información nueva se escribe en el cell state.
        \item Decidir cuánto del cell state actual se expone como salida $h_t$.
        \item Decidir qué información del cell state anterior se descarta (multiplicando por un valor entre 0 y 1).
        \item Calcular el gradiente para el backpropagation through time.
    \end{enumerate}
\end{enumerate}

\vspace{0.5cm}
\hrule
\vspace{0.5cm}

%% =============================================
\section*{Parte 3: Desarrollo (10 pts c/u)}

\begin{enumerate}[label=\arabic*)]
    \item \textbf{Forward Pass manual.} \\
    Considere una red con 2 entradas $(x_1, x_2)$, una capa oculta de 2 neuronas con activación ReLU, y una neurona de salida con activación Sigmoide. Los pesos y biases son:

    \begin{center}
    \begin{tabular}{lcc}
    \toprule
    & Neurona oculta 1 & Neurona oculta 2 \\
    \midrule
    $w_{x_1}$ & 0.5 & $-1.0$ \\
    $w_{x_2}$ & 1.0 & 0.5 \\
    $b$ & $-0.5$ & 0.0 \\
    \bottomrule
    \end{tabular}
    \end{center}

    Neurona de salida: $w_1 = 1.0$, $w_2 = -0.5$, $b = 0.2$.

    Para la entrada $x_1 = 1, x_2 = 0$:
    \begin{enumerate}[label=\alph*)]
        \item Calcule $z$ y $a$ para cada neurona de la capa oculta.
        \item Calcule $z$ y $a$ para la neurona de salida.
        \item Si el label verdadero es $y = 1$, ¿la red predice correctamente (umbral 0.5)?
    \end{enumerate}

    \item \textbf{Segmentación semántica y Dice Loss.} \\
    \begin{enumerate}[label=\alph*)]
        \item ¿Por qué en segmentación semántica se prefiere Dice Loss sobre Cross-Entropy clásica?
        \item Dada una predicción y un ground truth para una imagen de $4 \times 4$ píxeles (clase ``objeto'' = 1, ``fondo'' = 0):

        Predicción: \texttt{[[1,1,0,0],[1,1,0,0],[0,0,0,0],[0,0,0,0]]}

        Ground Truth: \texttt{[[1,1,1,0],[1,1,0,0],[1,0,0,0],[0,0,0,0]]}

        Calcule el coeficiente Dice: $\text{Dice} = \frac{2|X \cap Y|}{|X| + |Y|}$.
    \end{enumerate}

    \item \textbf{Transformers: Self-Attention.} \\
    Considere la oración: ``El banco está cerca del río''.
    \begin{enumerate}[label=\alph*)]
        \item Explique intuitivamente qué representan las matrices Query, Key y Value en el mecanismo de Self-Attention.
        \item ¿Cómo resuelve Self-Attention la ambigüedad de la palabra ``banco'' (¿institución financiera o ribera?) en esta oración?
        \item ¿Para qué sirve Multi-Head Attention? Dé un ejemplo de dos ``relaciones'' distintas que podrían aprender dos heads diferentes.
    \end{enumerate}
\end{enumerate}

\vspace{0.5cm}
\hrule
\vspace{0.5cm}

%% =============================================
\section*{Parte 4: Interpretación de Código (15 pts)}

Analice el siguiente código de PyTorch y responda las preguntas.

\begin{lstlisting}
import torch
import torch.nn as nn

class CNN(nn.Module):
    def __init__(self):
        super().__init__()
        self.conv1 = nn.Conv2d(1, 32, kernel_size=9, padding=4)
        self.conv2 = nn.Conv2d(32, 64, kernel_size=3, padding=1)
        self.pool = nn.MaxPool2d(2, 2)
        self.relu = nn.ReLU()
        self.fc1 = nn.Linear(64 * 7 * 7, 128)
        self.fc2 = nn.Linear(128, 10)

    def forward(self, x):
        x = self.pool(self.relu(self.conv1(x)))
        x = self.pool(self.relu(self.conv2(x)))
        x = x.view(-1, 64 * 7 * 7)
        x = self.relu(self.fc1(x))
        x = self.fc2(x)
        return x
\end{lstlisting}

\begin{enumerate}[label=\alph*)]
    \item ¿Cuál es el tamaño esperado de la imagen de entrada? Justifique rastreando las dimensiones capa por capa hasta llegar a $64 \times 7 \times 7$ antes del flatten. (Use la fórmula $O = \frac{I - K + 2P}{S} + 1$.)

    \item ¿Por qué el \texttt{padding=4} en la primera convolución y \texttt{padding=1} en la segunda? ¿Qué relación tiene con el tamaño del kernel?

    \item ¿Cuántos parámetros entrenables tiene la capa \texttt{conv1}? Muestre el cálculo.

    \item ¿Por qué la capa \texttt{fc1} tiene como entrada $64 \times 7 \times 7 = 3136$ neuronas?

    \item ¿Qué problema está resolviendo esta red? ¿Para qué dataset fue diseñada?

    \item Si quisiera adaptar esta red para clasificar imágenes a color de $32 \times 32$ en 100 clases (ej. CIFAR-100), ¿qué cambios mínimos haría en la arquitectura?
\end{enumerate}

\end{document}
